\input{preambel} 	% Alle Packages und Formatierungen werden in preamble.tex eingetragen um den Text hier übersichtlicher zu gestalten.
 
\section{Einleitung} % Überschrift für erstes Kapitel, optional auch: [Eintrag im IHV]{Einleitung}
\label{sec:Einleitung} % Textmarke, um darauf zu verweisen.
\lipsum[2]			% Dummytext (einfach löschen und los schreiben)

\section{Thematik in der Literatur}
\lipsum[30]				% Dummytext (einfach löschen und los schreiben)

Es gibt eine Reihe von Möglichkeiten chemische Formeln wie in \textsl{\ref{XXX}} zu notieren. Eine davon ist diese:

\begin{equation}					% Gleichung
\begin{aligned}						% Nummerierung am Rand
\label{XXX}							% beliebiges Label
\textnormal{XXX}_2 \hspace{0.3cm}\autorightarrow{drüber}{drunter}\hspace{0.3cm} \textnormal{XX}_2 + \textnormal{X}
\end{aligned}
\end{equation}
\vspace{0.2cm}

Wenn nur einzelne Ionen notiert werden sollen können diese entweder direkt im Text als $\anion{--II}{O}{}{2}$ geschrieben, oder in der Reaktionsgleichung \textsl{\ref{YYY}} notiert werden:

\begin{equation}					% Gleichung
\begin{aligned}						% Nummerierung am Rand
\label{YYY}							% beliebiges Label
\textnormal{O}_2 + 4\, \ee{}{} \hspace{0.3cm}\autorightarrow{Reduktion}{}\hspace{0.3cm} 2\,\anion{--II}{O}{}{2} 
\end{aligned}
\end{equation}
\vspace{0.2cm}

\begin{figure}[!ht]		% Bild mit Unterschrift einbinden
\centering
\includegraphics[scale=1]{pics/bild.png} 
\caption[Demobild]{Einzelne Bilder wie dieses sollten zentriert auf der Seite platziert werden. Auf eine sinnvolle Größe ist ebenso zu achten, wie auf eine konkrete Beschreibung der Abbildung.}
\label{fig:bild1}\vspace{0.8em}
\end{figure}

\lipsum[31]

\newpage
\section{Überschrift}

\renewcommand{\baselinestretch}{1.0}\normalsize 
\setlength{\columnsep}{30px} % Abstand Spalten vergrößern
\begin{multicols}{2}[\noindent \textbf{Hier könnte beispielsweise  ein Verweis auf den Textauszug aus der Literatur stehen.\citeup{voss}{}}] %\textbf{Dies kann als Zusammenfassung für den darunter liegenden Text gesehen werden. Hin und wieder ist dieser auch fett gedruckt.}
\setcounter{linenumber}{1}
\begin{linenumbers} % Schaltet Zeilennummerierung ein
%\modulolinenumbers[5] % nur jede 5. Zeile
\lipsum[99]
\end{linenumbers}
\end{multicols} % Ende 2 Spalten
\renewcommand{\baselinestretch}{1.5}\normalsize 

\subsection{Überschrift}
\lipsum[99]

\begin{table}[!ht]
\centering
\caption[Zahlenvergleich der Werte 1 und 2\citeup{voss}{}]{Die Werte sind an fiktiven Proben gemessen.}\vspace{0.8em}
\label{tab:tab1}
\begin{tabular}{l|cc} \toprule
 &\multicolumn{2}{r}{universal}\\ \cmidrule{1-1}
\backslashbox{ABC}{DEF} & 1 & 2 \\ \midrule 
E & 3,34 & 4,34 \\
F & 5,45 & 6,45 \\ \bottomrule
\end{tabular}
\end{table}

\lipsum[99]

\newpage
\section{Ausblick und Zusammenfassung}
Dies ist eine Text\footnote{Dies solle Fußnote werden, die den Rahmen einer Zeile sprengt. Dies lässt sich auch durch Wiederholung den Inhalts erreichen. Dies lässt sich auch durch Wiederholung den Inhalts erreichen.  Dies lässt sich auch durch Wiederholung den Inhalts erreichen. \newline Dies lässt sich auch durch Wiederholung den Inhalts erreichen.}, der auch in der Literatur existiert.\citeup{demo}{,S.22} 
\lipsum[32]

Da pack ich doch gleich no ne Fußnote drunter.\footnote{Mal sehen, welcher Abstand hier zu vorherigen Fußnote entsteht.}

\lipsum[10]

\begin{table}[!ht]					%	Tabelle genau hier einsetzen
\centering							% Tabelle zentrieren
\caption[YYY]{Tabellenüberschrift.}\vspace{0.8em} % Tabellenbeschrift.
\label{tab:2}							% Label für den Tabellenverweis
\begin{tabular}{l|r|r} \toprule		% l = linksb., c = zentriert, r = rechtsb., | vert. Strich
  & \textbf{Spalte A} & \textbf{Spalte B} \\ \midrule 				% & ist Spaltentrennzeichen, Abschluss der Zeile mit "\\"
\textbf{Zeile 1} & Wert 1A & Wert 1B \\
\textbf{Zeile 2 }& Wert 2A & Wert 2B \\ \bottomrule		% top-, mid- & bottomrule in untersch. Strichstärken
\end{tabular}
\end{table}

\lipsum[41]


\newpage
\section{Selbständigkeitserklärung}
Hiermit erkläre ich, dass ich die vorliegende Hausarbeit selbständig verfasst und keine anderen als die angegebenen Hilfsmittel benutzt habe.\footnote{eine kleine Fußnote.}

Die Stellen der Hausarbeit, die anderen Quellen im Wortlaut oder dem Sinn nach entnommen wurden, sind durch Angaben der Herkunft kenntlich gemacht. Dies gilt auch für Zeichnungen, Skizzen, bildliche Darstellungen sowie für Quellen aus dem Internet.

\vspace{20mm}
\hrule
\vspace{5mm}
\begin{center}
\textbf{Bad Reichenhall}, \today
\end{center}


\newpage
\listoffigures		% Abbildungsverzeichnis erstellen

\newpage
\listoftables		% Tabellenverzeichnis erstellen

\newpage
\addcontentsline{toc}{section}{Literaturverzeichnis}
\small 						%\documentclass[journal = jacsat]{achemso}
							%\setkeys{acs}{articletitle = true}
\bibliographystyle{angewchem} %achemso angewchem 
\bibliography{standard}

\newpage
\appendix
\addcontentsline{toc}{section}{Anhang}
\begin{appendix}
\section*{Anhang}
%\subsection{Anhangsteil 1}
%\lipsum[31]
%\subsection{Anhangsteil 2}
\end{appendix}
\end{document}